\clearpage

\section{二次型}

\subsection{从代数到几何}

目前我们已经通过线性空间的概念刻画了“向量”这一代数对象.
但是这一刻画在应用上还是有很大的不足.
回想中学时，
我们用向量来做解析几何问题，
经常要计算向量的长度、夹角等几何量.
但这些在线性空间中都还没有定义.
我们已经看到，
线性空间的概念是很泛用的，
函数、集合、量子态等等类型的元素都可以建立线性空间——
如果它们也能像解析几何那样谈论“长度”、“夹角”等概念，
岂不是很方便？

那么，在抽象的线性空间中，
我们没法直接看到“长度”、“夹角”这些东西，
如何去定义呢？
就像我们定义线性空间中的向量那样，
我们并不是直接去问“向量是什么”，
而是去问“向量满足哪些性质”.
这是代数学中非常重要的思想方法.

在欧氏空间中，
\textbf{内积}可以诱导出我们想要的这些概念：
一个向量与自身的内积即是其长度的平方；
两个向量的夹角余弦也可用内积来计算.
这样看来，与“距离”、“长度”相比，
“内积”是“级别更高”的一种结构.
为了把所有这些概念一网打尽，
我们直接在线性空间的基础上去定义\textbf{内积}.
定义了内积的线性空间称为\textbf{内积空间}.

\subsubsection{实内积空间概述}

内积结构的定义与数域有关.
一般我们只关注两个数域：$\mathbb{R}$和$\mathbb{C}$.
因此，
我们需要分别定义实线性空间和复线性空间的内积.
在线性代数中，
我们主要讨论实内积和实内积空间.
我们将在本节的最后再讨论$\mathbb{R}$到$\mathbb{C}$的推广.

\bigskip


以我们最熟悉的$\mathbb{R}^n$中的标准内积
（两个数组对应分量相乘再全部相加）为模板，
可以总结出内积“应有”的几条性质.
我们借助这些性质去叙述一般实内积的定义：

\medskip

\textit{设$V$是$\mathbb{R}$上的线性空间.
定义$V$上的一个二元函数$(\cdot,\cdot) : V \times V \to \mathbb{R}$，即：
对于任意$\bm{x},\bm{y} \in V$，$(\bm{x},\bm{y})$都对应一个实数.
如果对任意的$\bm{x},\bm{y} \in V$，$\alpha \in \mathbb{R}$，
满足以下$4$条规则，
那么称$V$是实内积空间（欧氏空间），
这个二元函数是$V$上的一个实内积：}

(1)~\textbf{正定性：}\textit{$(\bm{x},\bm{x} )\ge 0$，等号成立当且仅当$\bm{x} = \bm{0}$.}

(2)~\textbf{对称性：}\textit{$(\bm{x},\bm{y}) = (\bm{y},\bm{x})$.}

(3)~\textbf{线性（数乘）：}\textit{$(\alpha \bm{x}, \bm{y}) = \alpha (\bm{x},\bm{y})$.}

(4)~\textbf{线性（加法）：}\textit{$(\bm{x}+\bm{y},\bm{z}) = (\bm{x},\bm{z}) + (\bm{y},\bm{z})$.}\vspace{1em}

\medskip

$\mathbb{R}^n$中的标准内积是大家所熟悉的.
这里再举两个大家所知道的，
但可能未“如此想过”的内积空间.

\medskip

\textbf{例1}~~
\textit{全体$m \times n$实矩阵构成$\mathbb{R}$上的线性空间$M_{m \times n}(\mathbb{R})$.
在其上定义二元函数$(A,B) = \mathrm{tr}(A^T B)$. 
这就是$M_{m \times n}(\mathbb{R})$上的一个内积，称为Frobenius内积.}

\textbf{例2}~~
\textit{定义在$[a,b]$上的全体平方可积函数$L^2[a,b]$构成$\mathbb{R}$上的线性空间.
对$f,g\in L^2[a,b]$，定义：
}
\begin{align*}
    (f,g) = \int_{a}^{b} f(x) g(x) ~\mathrm{d}x
\end{align*}

\textit{则不难验证这是一个内积.}

上面的例2来自于高等数学，
以此为基础建立了傅里叶级数理论.


\medskip

在定义了内积之后，
便可以获得许多我们在$\mathbb{R}^n$中所熟悉的概念：

\textbf{长度（范数）}
~~
\textit{设$V$是实内积空间，$\bm{\alpha} \in V$，
则$\bm{\alpha}$的长度（范数）定义为：$\left\| \bm{\alpha} \right\| = (\bm{\alpha} , \bm{\alpha})^{1/2}$.}

\textbf{距离（度量）}
~~
\textit{设$V$是实内积空间，$\bm{\alpha},\bm{\beta} \in V$，
则$\bm{\alpha}$与$\bm{\beta}$的距离定义为：$d (\bm{\alpha} , \bm{\beta}) 
= \left\| \bm{\alpha} - \bm{\beta} \right\|$.}

\smallskip

\textbf{夹角}
~~
\textit{设$V$是实内积空间，
定义非零向量$\bm{\alpha}$和$\bm{\beta}$的夹角$\theta$之余弦为$\cos \theta = \frac{(\bm{\alpha}, \bm{\beta})}{\left\| \bm{\alpha}\right\| \left\| \bm{\beta} \right\|}$.}


\medskip


\textbf{定理4.4.1}~~
\textit{设$V$是实内积空间，
$\bm{\alpha},\bm{\beta} \in V$，
$c$是常数，则：}

\textit{(1)~~$\left\| c \bm{\alpha}\right\| = |c| \left\| \bm{\alpha}\right\|$}.

\textit{(2)~~$|(\bm{\alpha} , \bm{\beta})| \leq \left\| \bm{\alpha} \right\| \cdot \left\| \bm{\beta} \right\|$ }.

\textit{(3)~~$\left\| \bm{\alpha} + \bm{\beta} \right\| \leq \left\| \bm{\alpha} \right\| + \left\| \bm{\beta} \right\|$}.

\subsubsection{内积的表示——正定实对称矩阵}

对于一般的$n$维欧氏空间或酉空间，
在给定一组基$\{ \bm{e}_1 , \bm{e}_2 , \cdots , \bm{e}_n \}$之后，
其中的向量就能表示为坐标（数组）.
于是我们希望能利用坐标去计算内积.

以下设$\bm{\alpha} = x_1 \bm{e}_1 + x_2 \bm{e}_2 + \cdots + x_n \bm{e}_n$，
$\bm{\beta} = y_1 \bm{e}_1 + y_2 \bm{e}_2 + \cdots + y_n \bm{e}_n$.
它们在基$\{ \bm{e}_1 , \bm{e}_2 , \cdots , \bm{e}_n \}$下的坐标分别为$\bm{X}$和$\bm{Y}$.
记$(\bm{e}_i , \bm{e}_j) = g_{ij}$.
因为实内积有线性性质，
所以任意两个向量$\bm{\alpha},\bm{\beta}$
的内积可以由一组基两两之间的内积$(\bm{e}_i , \bm{e}_j) = g_{ij}$
来确定：
\begin{align*}
    (\bm{\alpha} , \bm{\beta}) = 
    \left( \sum_{i=1}^{n} x_i \bm{e}_i , \sum_{j=1}^{n} y_j \bm{e}_j \right)
    = \sum_{i=1}^{n} \sum_{j=1}^{n} g_{ij} x_i y_j 
\end{align*}

这个结论实际上可以写成矩阵形式：
\begin{align*}
    (\bm{\alpha} , \bm{\beta})
    = (x_1 , x_2 , \cdots ,x_n) 
    \begin{pmatrix}
        g_{11} & g_{12} & \cdots & g_{1n} \\
        g_{21} & g_{22} & \cdots & g_{2n} \\
        \vdots & \vdots  & \ddots & \vdots \\
        g_{n1} & g_{n2} & \cdots & g_{nn}
    \end{pmatrix}
    \begin{pmatrix}
        y_1 \\
        y_2 \\
        \vdots \\
        y_n
    \end{pmatrix}
\end{align*}

其中矩阵$G$：
\begin{align*}
    G = 
    \begin{pmatrix}
        g_{11} & g_{12} & \cdots & g_{1n} \\
        g_{21} & g_{22} & \cdots & g_{2n} \\
        \vdots & \vdots  & \ddots & \vdots \\
        g_{n1} & g_{n2} & \cdots & g_{nn}
    \end{pmatrix}
    =
    \begin{pmatrix*}
        (\bm{e}_1, \bm{e}_1)  &  (\bm{e}_1, \bm{e}_2) & \cdots &  (\bm{e}_1, \bm{e}_n) \\
        (\bm{e}_2, \bm{e}_1) &  (\bm{e}_2, \bm{e}_2) & \cdots &  (\bm{e}_2, \bm{e}_n) \\
        \vdots & \vdots  & \ddots & \vdots \\
         (\bm{e}_n, \bm{e}_1) &  (\bm{e}_n, \bm{e}_2) & \cdots &  (\bm{e}_n, \bm{e}_n)
    \end{pmatrix*}
\end{align*}

这里的矩阵$G$也称为\textbf{度规矩阵}.
$\bm{\alpha}, \bm{\beta}$在给定基下的内积即可表示为$(\bm{\alpha} , \bm{\beta}) = \bm{X}^T G \bm{Y}$.
容易看出$G$是对称矩阵.
另外由内积的正定性可知，
$\bm{X}^T G \bm{X} > 0$对于任意非零向量$\bm{X}$成立.
满足这种性质的矩阵$G$称为\textbf{正定矩阵}.
表达式$\bm{X}^T G \bm{X}$称为\textbf{二次型}，
它是关于$\bm{X}$各分量的二次齐次多项式.

由此看出，
\textbf{欧氏空间中的实内积与正定实对称矩阵之间存在一一对应.}


\subsubsection{广义内积的引入与表示——实对称矩阵}

在我们推广内积的概念之前，
我们先对\textbf{度规矩阵}$G$的意义做进一步讨论.
我们知道$G$的$(i,j)$元素是$(\bm{e}_i , \bm{e}_j)$.
根据长度和夹角的定义，
$G$对角线上的元素是定义了一组基中每个$\bm{e}_i$的长度，
而其他位置的元素是定义了这组基两两之间的夹角.
因此，
$G$其实打包了一组基的全部几何信息，
也就代表了整个空间的几何信息.
它的作用就像是一把测量标尺.
这就是为什么我们称其为度规矩阵.

度规矩阵的另一种（本质相同的）看法是给出了内积空间中的基本几何定律.
比如在我们所熟悉的欧氏空间中，
建立解析几何的基本定律是勾股定理.
这与欧氏度规（单位矩阵$I$）是等价的.
勾股定理可以表达为：
\begin{align*}
    \left\| \bm{X} \right\|^2 = \bm{X}^{T} I \bm{X}
\end{align*}

勾股定理的一个重要应用是计算曲线长.
我们把曲线微元$\mathrm{d} s$近似为直线段，
则能得出线元$\mathrm{d} s^2 = (\mathrm{d} x)^2 + (\mathrm{d} y)^2$.
这在高等数学中有详细说明.

如果我们把欧氏度规$I$换成任意的度规$G$，
我们会得到相应内积空间中的“勾股定理”
\begin{align*}
    \left\| \bm{X} \right\|^2 = \bm{X}^{T} G \bm{X}
\end{align*}


在高等数学中，
我们可以讨论如何在球坐标下计算三维曲线的长度.
这个问题其实就是取球面度规：
\[
B = 
\begin{pmatrix}
    1 & 0 & 0 \\
    0 & r^2 & 0 \\
    0 & 0 & r^2 \sin^2 \theta
\end{pmatrix}
\]

那么曲线的线元公式就可以写成：
\[
\mathrm{d} s^2 = 
(\mathrm{d} r , \mathrm{d} \theta, \mathrm{d} \varphi)
\begin{pmatrix}
    1 & 0 & 0 \\
    0 & r^2 & 0 \\
    0 & 0 & r^2 \sin^2 \theta
\end{pmatrix}
\begin{pmatrix}
    \mathrm{d} r \\
    \mathrm{d} \theta \\
    \mathrm{d} \varphi \\
\end{pmatrix}
= 
(\mathrm{d}r)^2 + r^2 (\mathrm{d} \theta)^2 + r^2 \sin^2 \theta (\mathrm{d} \varphi)^2
\]

有时我们也把线元称为\textbf{度规}，
即通过线元可以直接“读”出度规矩阵，
从而知道当前坐标系下的几何特性.
给出球面度规即给出了球面几何.
我们容易想象在球面上的几何与通常的平面几何（欧氏几何）有许多不同之处（例如平行线会相交），
而度规则给了我们计算非欧几何的基本工具.

到目前为止，
我们所讨论的度规矩阵都是正定的.
这是因为实内积本身是正定的.
在我们通常的观念里，
虽然有球面几何这样“弯曲”的几何学，
但是“距离的平方是非负数”总归是毫无争议的.
然而随着科学的发展，
我们渐渐发现，
不满足正定性的几何学竟然也是有用的.
最经典的例子就是狭义相对论中，
“时空距离”并不符合$\mathbb{R}^4$中的勾股定理，
而是：
\begin{align*}
    \d s ^2 = -(\mathrm{d}t)^2 + ( \mathrm{d} x)^2 + ( \mathrm{d} y)^2 + (\mathrm{d} z)^2
\end{align*}

这里我们取光速$c =1$.
从这个线元中“读”出的度规是\textbf{闵氏度规}：
\[\eta = 
\begin{pmatrix}
    -1 & 0 & 0 & 0 \\
    0 & 1 & 0 & 0 \\
    0 & 0 & 1 & 0 \\
    0 & 0 & 0 & 1 
\end{pmatrix}
\]

按照定义它显然不是正定矩阵.
广义相对论进一步认为，
引力本质上是时空的“弯曲”，
即物质通过改变时空的几何（度规）来导致运动.
这使得非欧几何，而且是非正定的几何，
称为描述相对论物理的基本语言.

在这一背景下，
我们有必要对实内积的概念推广为\textbf{广义内积}，
或称为\textbf{对称双线性形式}.
推广的方式也很简单：在定义中去掉正定性即可.
相应地，
其表示矩阵一般也不具有正定性的要求，
而是一般的实对称矩阵$A$.
在$A$所规定的几何中，
“勾股定理”表达为$\left\| \bm{X} \right\|^2 = \bm{X}^T A \bm{X}$.
其中，
等号右边即是“斜边长的平方”的展开式，
这种表达式称为\textbf{二次型}，
$A$也可称为\textbf{二次型矩阵}.
如果写出分量式$\bm{X} = (x_1 , x_2 , \cdots , x_n)^T$，
那么二次型可以看成是一个$n$元的$2$次齐次多项式：
\begin{align*}
    f(x_1 , x_2 , \cdots , x_n)
    = \sum_{i=1}^{n} \sum_{j=1}^{n} a_{ij}x_i x_j
\end{align*}

线元$\mathrm{d} s^2$也可以看成一个二次型，
只不过代入了向量微分$\mathrm{d} \bm{X} = (\mathrm{d} x_1 , \mathrm{d} x_2 , \cdots , \mathrm{d} x_n)$.
更一般地，
如下形状的二次微分式也可以当作二次型来研究：
\begin{align*}
    f(\mathrm{d} x_1 , \mathrm{d} x_2 , \cdots , \mathrm{d} x_n)
    = \sum_{i=1}^{n} \sum_{j=1}^{n} a_{ij}\mathrm{d} x_i \mathrm{d} x_j
\end{align*}





\subsubsection{*自伴算子}

矩阵是对线性映射的表示.
但现在我们又说，
广义实内积可以表示为实对称矩阵.
两种看法能不能统一起来呢？
考察广义实内积的矩阵表达式：

\begin{align*}
    (\bm{\alpha} , \bm{\beta})
    = \bm{X}^T A \overline{\bm{Y}}
    = 
    (x_1 , x_2 , \cdots , x_n) 
    \begin{pmatrix}
        a_{11} & a_{12} & \cdots & a_{1n} \\
        a_{21} & a_{22} & \cdots & a_{2n} \\
        \vdots & \vdots  & \ddots & \vdots \\
        a_{n1} & a_{n2} & \cdots & a_{nn}
    \end{pmatrix}
    \begin{pmatrix}
        y_1 \\
        y_2 \\
        \vdots \\
        y_n
    \end{pmatrix}
\end{align*}

我们可以这样看待这个表达式：
它无非是在一个标准内积式$\bm{X}^T \bm{Y}$
中间插入了一个实对称矩阵$A$，
变成$\bm{X}^T A \bm{Y}$.
如果把$A$看成欧氏空间$V$上的一个线性变换$\mathcal{A}$，
那么我们会获得一个关于内积的全新理解：
\textbf{欧氏空间中的任意广义内积结构都可以看成该空间上的一个线性变换$\mathcal{A}$，
两个向量$\bm{\alpha},\bm{\beta}$的广义内积就是让$\mathcal{A}$
作用在其中一个向量上，
再让二者做$\mathbb{C}^n$中的标准内积.}




$A$是对称矩阵这一点对$\mathcal{H}$提出了一个限制.
如果欧氏空间中的线性变换能表示为实对称矩阵，
那么称这种变换为\textbf{自伴算子}.
相应的表示矩阵（实对称矩阵）称为\textbf{自伴矩阵}.
注意这里实质上没有引入新的数学概念，
只是借助“自伴”的名称引入了一种对于实对称矩阵的新观点.
任意的广义内积可以看成是
某个\textbf{自伴算子加权下的标准内积}.
自伴算子的性质及其与内积的关系可由下式表示：
\begin{align*}
    (\bm{\alpha} , \bm{\beta})
    =
    (\mathcal{A}\bm{\alpha} , \bm{\beta})_{\mathbb{R}}
    =
    (\bm{\alpha} , \mathcal{A}\bm{\beta})_{\mathbb{R}}
\end{align*}

其中$(\cdot , \cdot)$是欧氏空间$V$上的某个内积，
它在给定基下表示为矩阵$A$.
而$(\cdot , \cdot)_{\mathbb{R}}$
是指两个向量在给定基下的坐标下进行$\mathbb{R}^n$上的标准内积.

我们借助矩阵表达式进一步地理解这一点.
按矩阵的乘法结合律，
表达式$\bm{X}^T A \bm{Y}$有两种结合方式：
\begin{align*}
    (x_1 , x_2 , \cdots ,x_n) 
    \boxed{
         \begin{pmatrix}
        a_{11} & a_{12} & \cdots & a_{1n} \\
        a_{21} & a_{22} & \cdots & a_{2n} \\
        \vdots & \vdots  & \ddots & \vdots \\
        a_{n1} & a_{n2} & \cdots & a_{nn}
    \end{pmatrix}
    \begin{pmatrix}
        y_1 \\
        y_2 \\
        \vdots \\
        y_n
    \end{pmatrix}
    }
    = \boxed{
        (x_1 , x_2 , \cdots ,x_n) 
         \begin{pmatrix}
        a_{11} & a_{12} & \cdots & a_{1n} \\
        a_{21} & a_{22} & \cdots & a_{2n} \\
        \vdots & \vdots  & \ddots & \vdots \\
        a_{n1} & a_{n2} & \cdots & a_{nn}
    \end{pmatrix}
    }
    \begin{pmatrix}
        y_1 \\
        y_2 \\
        \vdots \\
        y_n
    \end{pmatrix}
\end{align*}

左式为$\bm{X}^T (A \bm{Y})$，
即如果让$\mathcal{A}$作用在$\bm{\beta}$上（坐标为$\bm{Y}$），
则可以看出$\mathcal{A}$的表示矩阵应为$A$.
右式为$(\bm{X}^T A) \bm{Y} = (A^T \bm{X})^T \bm{Y}$.
即如果让$\mathcal{A}$作用在$\bm{\alpha}$上（坐标为$\bm{X}$），
则可以看出$\mathcal{A}$的表示矩阵应为$A^T$.
要求$A$是自伴矩阵，
即是要求$A = A^T$是$\mathcal{A}$的同一个表示矩阵.


\subsubsection{矩阵的合同}

我们知道（广义）内积的定义是无关于坐标系（基）的选取的，
但是度规矩阵的写法却与基的选择有关.
这与线性映射及其表示矩阵的关系类似，
但是在内积背景下，
就给出了一种不同于矩阵相似的基变换关系.

设$n$维欧氏空间中规定了一个（广义）内积$(\cdot , \cdot)$.
它在基$\{ \bm{e}_i \}$下的度规矩阵为$A$.
向量$\bm{\alpha}, \bm{\beta}$在这组基下的坐标分别为$\bm{X}$和$\bm{Y}$.
则$(\bm{\alpha} , \bm{\beta}) = \bm{X}^T A \bm{Y}$.
若基变换$\{ \bm{e}_i \} \mapsto \{ \bm{e}'_i \}$
以$P$为过渡矩阵，
新旧坐标之间的关系
$P\bm{X}' = \bm{X}, ~ P\bm{Y}' = \bm{Y}$.
于是有：
\begin{align*}
    (\bm{\alpha} , \bm{\beta})  = \bm{X}^T A \bm{Y}
    = (P\bm{X}')^T A P\bm{Y}'
    = \bm{X}'(P^T A P) \bm{Y}'
\end{align*}

可见同一个内积在新基下的度规矩阵为$P^T A P$.
也就是说，
对于实对称矩阵$A$，
如果存在$n$阶可逆矩阵$P$使得
$B = P^{T}AP$（此时$B$必然也是实对称的），
则矩阵$A$和$B$是
\textbf{同一个广义内积在不同基下的度规矩阵}.
此时我们称矩阵$A$和$B$是\textbf{合同}的.
变换$A \mapsto P^{T}AP$称为\textbf{合同变换}.
可以证明相似变换是一种等价关系：

\vspace{0.5em}


(1)~~反身性：取$P = I$，显然$I^{T} A I = A$.

(2)~~对称性：如果$A$合同于$B$，
即存在可逆矩阵$P$使得$B = P^{T}AP$.
取$Q = P^{-1}$，则$A = Q^{T}BQ$，
因此$B $与$A$合同.

(3)~~传递性：
如果$A$与$ B$合同,
$B$与$C$合同，
即存在可逆矩阵$P,Q$，
使得$B = P^{T} A P, ~C = Q^{T} B Q$，
则$C = Q^{T} P^{T} A PQ$.
取$R = PQ$，
则$C = R^{T} A R$，
即$A $与$C$合同.

\vspace{0.5em}

如果把实对称矩阵$A$考虑成一个二次型矩阵，
那么我们说两个\textbf{二次型等价}，
就等同于说它们的二次型矩阵是合同的.
当我们做坐标变换$\bm{X} = P\bm{X}'$时（$P$是可逆矩阵），
在二次型的视角看来，
其实就是做了一个变量替换
$(x_1 , x_2 , \cdots , x_n) \mapsto (x_1' , x_2' , \cdots , x_n')$.
这种用可逆矩阵定义的变量替换称为\textbf{非退化线性替换}.


度规矩阵在合同关系下的不变量称为\textbf{合同不变量}.
同一个广义内积在不同的基下对应着不同的度规矩阵，
但是这些度规矩阵应该有共同的，
与基的选择无关的性质.
度规矩阵的这类性质可以看成是（广义）内积结构本身的属性，
这就是合同不变量的意义.
关于“不变量”思想的理解，
请与相似不变量相互参考.
好消息是，
合同不变量比相似不变量简单得多，
可以找出完全不变量.
也就是说，
我们在线性代数的课程范围内就可以做到对矩阵在合同意义下进行完全分类，
并找出所谓“合同标准型”.
这等价于对欧氏空间中所有可能的（广义）内积结构进行了完全分类.
我们将在下一节详细讨论.


\subsubsection{保度规变换与正交矩阵}

欧氏空间中的线性变换可能会改变空间的几何.
例如在$\mathbb{R}^2$中直观地看，
伸缩和剪切会导致形变，
而反射和旋转却不会.
类似于后者的变换是特别重要的.
例如物理中的刚体运动，
正是保持物体内部的长度，夹角等所有几何性质完全不变的一种运动.
我们提出\textbf{保内积}
的概念来描述这种变换.
我们将这个性质抽象成如下的数学描述：

\medskip

\textit{如果欧氏空间上规定了（广义）内积$(\cdot , \cdot)$，
在给定基下的度规矩阵为$A$.
存在线性变换$\mathcal{M}$，
满足对任意的$\bm{\alpha}, \bm{\beta}$（坐标分别为$\bm{X},\bm{Y}$），
有：}

\[
(\mathcal{M}\bm{\alpha} , \mathcal{M}\bm{\beta}) = (\bm{\alpha} , \bm{\beta})
\]

\textit{那么$\mathcal{M}$称为保内积变换}.

\medskip

在这个话题中，
我们总是在用一组基下描述问题，不涉及基的变换.
因此所谓“保内积”就是“保度规”.
设$\mathcal{M}$的表示矩阵为$M$，
则等价地有$M$满足：

\[
(M\bm{X})^T A (M\bm{Y}) = \bm{X}^T A \bm{Y},
~~~~~\text{即}
\bm{X}^T M^T A M \bm{Y}= \bm{X}^T A \bm{Y}
\]

由此可以看出$M$应该满足$M^T A M = A$，
这就是“保度规”的意思.

\medskip

作为具体例子，
我们考察欧氏度规（单位矩阵）$I$的保度规变换.

根据定义，它满足$M^T I M = I$，
即$M^T = M^{-1}$.
有这一性质的变换$\mathcal{M}$称为\textbf{正交变换}，
表示矩阵$M$称为\textbf{正交矩阵}.
所谓正交变换，
就是指保持欧氏距离和夹角不变的变换，
简而言之就是旋转或者反射变换.
因此，不难想象正交矩阵应该长什么样.
我们说，
一个线性变换可以由它在一组标准正交基上的作用唯一确定，
表现为矩阵的列向量组.
而正交矩阵代表了保持欧氏距离和夹角不变的变换.
因此，正交矩阵的列向量应该是一组单位正交基，
即每一个列向量的欧氏长度为$1$，
且两两正交.

不出意外的话，
欧氏几何是大家唯一熟悉的几何，
因此正交变换也是大家唯一能够直观理解的保度规变换.
但是对于学过狭义相对论的同学，
不妨思考保度规变换在其中的应用：
狭义相对论的时空几何是由闵氏度规$\eta = \mathrm{diag}(-1,1,1,1)$来描述.
保闵氏度规的变换$\Lambda$满足$\Lambda^T \eta \Lambda = \eta$.
$\Lambda$就是\textbf{洛伦兹变换}.
洛伦兹变换给出了不同惯性系下时空坐标变换的规律.
于是我们知道，
惯性系的切换只是一种$4$维时空的“转动”.
在不同的惯性系下，
（闵氏度规下的）时空距离是不变量.


\subsection{矩阵合同的进一步讨论}

\subsubsection{合同标准型和规范型}

设$A$是$n$阶实对称矩阵.
为了研究矩阵在合同意义下的分类，
与上一章类似地，
我们提出课题：
能否找到一个可逆矩阵$C$，
使得$C^T A C$是对角矩阵？
与之等价的问题是，
对于一个二次型$\bm{X}^T A \bm{X}$，
能否找到一个非退化线性替换$\bm{X} = C \bm{Y}$，
使得替换后的二次型只含有平方项而不含有交叉项？
我们把这种对角矩阵（只有平方项的二次型）
称为矩阵$A$（或者二次型$\bm{X}^T A \bm{X}$）的\textbf{合同标准型}.

任何实对称矩阵都能化成合同标准型吗？
答案是可以，
而且不难给出理由：
我们在上一章就知道，
实对称矩阵一定可以正交对角化.
所谓正交对角化是指，
存在正交矩阵$C$，
使得$C^{-1} A C$为对角矩阵.
但是注意正交矩阵的性质$C^{T} = C^{-1}$，
$C^{-1} A C$其实就是$C^{T} A C$.
也就是说，
\textbf{把实对称矩阵的正交对角化，其实就是化为了合同标准型.}
\textbf{实对称矩阵的相似标准型就是一个合同标准型.}
因此，
计算合同标准型根本不需要新的方法，
只需要按照上一章的方法，
去计算正交对角化即可.
当然，也有别的办法更直接地去操作合同变换，
进而得到合同标准型，
将在习题部分介绍.

这里需要澄清的是，
\textbf{实对称矩阵的合同标准型并不是唯一的}，
可逆矩阵$C$的选取也不是唯一的.
因此，
虽然实对称矩阵总是与其相似标准型是合同的，
也并不意味着相似不变量一定就是合同不变量，
只能说相似不变量可以反映某些合同的性质.
具体将在下一小节中分析.

\medskip

我们已经能找到实对称矩阵$A$的一个合同标准型$D$.
设$D$是如下的对角矩阵：
\begin{align*}
    D = \mathrm{diag}(d_1 , \cdots , d_p , -d_{p-1} , \cdots , -d_r , 0 , 0 \cdots , 0)
\end{align*}

对应的二次型是：
\begin{align*}
    f = d_1 x_1 ^2 + \cdots + d_p x_p^2 - d_{p+1}x_{p+1}^2 - \cdots - d_{r} x_{r}^2
\end{align*}

方便起见，这里我们已将对角线上的正数、负数和零分别排列到一起（这总是可以做到的）.
也就是说这里的$d_i$全都为正数，$r \leq n$.
不难想到，
可以对二次型做进一步的化简：
\begin{align*}
    x_i = 
    \begin{cases}
        \frac{y_i}{\sqrt{d_i}} , \quad 1 \leq i \leq r \\
        y_i ,\quad r < i \leq n
    \end{cases}    
\end{align*}

得到的对角矩阵，对角线上只含$1,-1$和$0$：
\begin{align*}
    D = \mathrm{diag}(1 ,  \cdots , 1 , -1 , \cdots , -1 , 0 , 0 \cdots , 0)
\end{align*}

对应的二次型是：
\begin{align*}
    f = y_1 ^2 + \cdots + y_p^2 - y_{p+1}^2 - \cdots - y_{r}^2
\end{align*}

这似乎是我们能得到形式最简单的东西.
我们将上述形状的对角矩阵（二次型）称为实对称矩阵$A$（二次型$\bm{X}^T A \bm{X}$）的\textbf{规范型}.
从直观上，
标准型之所以不唯一，
是因为系数的“大小”可以是任意的.
究其根本，
是说对于同一个（广义）内积结构而言，
用度规矩阵“测量出来的长度”并不是绝对的，
而是可以通过选取不同的基来任意确定.
这就好比一根1 m长的木棍，
同时也是10 dm和100 cm——
不同单位的选取会让长度的“读数”发生变化.
因此，
对角元的绝对大小在合同关系中并不重要.
把标准型进一步变换为规范型，
正是剔除了这一冗余的信息.

一个$n$阶实对称矩阵的规范型，
其实只需要两个数字就能确定：
$p$和$r$.
其中$r$是非零元（$\pm 1$）的个数，
其实也就是秩（注意秩也是合同不变量）.
$p$则是$1$的个数，
我们称为\textbf{正惯性指数}.
方便起见，
我们还约定规范型中$-1$的个数$r-p$为\textbf{负惯性指数}，
正惯性指数减去负惯性指数所得的差$2p-r$称为\textbf{号差}.

\subsubsection{规范型的进一步讨论}


在教材中证明了，
\textbf{$n$阶实对称矩阵的规范型总是存在且唯一的.}
也就是说，
谈论一个实对称矩阵作为（广义）内积的性质，
就等同于谈论其规范型的性质.
根据上一小节最后的讨论，
\textbf{秩}与\textbf{正惯性指数}组成了二次型（矩阵）的一组完全不变量.
我们可以分别从主动观点和被动观点来理解这一点.
从主动观点来看：
\textbf{不同的（广义）内积结构，
本质上只有秩和正惯性指数的不同.}
从被动观点来看：
\textbf{同一个内积结构虽然对应不同的度规矩阵，
但是它们的秩和正惯性指数是确定的.}

这里我们以闵氏度规$\eta = \mathrm{diag}(-1,1,1,1)$为例简要说明这个结论的意义.
在闵氏度规下，
三个空间轴对应的$4$维矢量
$(0,x,0,0)^T,(0,0,y,0)^T,(0,0,0,z)^T$长度的平方为正数，
这反映了$3$维空间具有我们所熟悉的实内积性质.
而时间轴对应的$4$维矢量$(t,0,0,0)^T$长度的平方为负数.
借用这些物理上的背景，
我们提出以下概念：
在给定了广义内积结构的实线性空间中，
长度的平方为正的矢量称为\textbf{类空矢量}；
长度的平方为负的矢量称为\textbf{类时矢量}；
长度的平方为零的矢量称为\textbf{类光矢量}.
\textbf{一个广义内积结构最本质的作用是将整个空间划分为类空区域、类时区域和类光区域.}
我们可以不太严谨地说，
所谓正/负惯性指数，
就是指整个空间中类空/类时区域的维数.
比如，
从闵氏度规中可以读出，
闵氏时空中有$1$个时间维度，$3$个空间维度.

以上我们明确了，可以使用规范型作为合同等价类的代表元.
我们引入以下的词汇来描述不同类型的实对称矩阵（广义内积结构/规范型）：

\medskip

(1)~实对称矩阵$A$是\textbf{正定}的，当且仅当以$A$为度规的广义内积是正定的（即：$A$是一个实内积的度规矩阵），
当且仅当$A$的规范型的对角元全为$1$（单位矩阵$I_n$）.

(2)~实对称矩阵$A$是\textbf{负定}的，当且仅当以$A$为度规的广义内积是负定的（非零向量长度的平方为负数），
当且仅当$A$的规范型的对角元全为$-1$（负单位矩阵$-I_n$）.

(3)~实对称矩阵$A$是\textbf{半正定}的，当且仅当以$A$为度规的广义内积是半正定的（非零向量长度的平方为非负数），
当且仅当$A$的规范型的对角元全为$1$和$0$.

(4)~实对称矩阵$A$是\textbf{半负定}的，当且仅当以$A$为度规的广义内积是半负定的（非零向量长度的平方为非正数），
当且仅当$A$的规范型的对角元全为$-1$和$0$.

(5)~实对称矩阵$A$是\textbf{不定}的，当且仅当以$A$为度规的广义内积是不定的（非零向量长度的平方可正可负），
当且仅当$A$的规范型的对角元既有$1$又有$-1$.
不定矩阵可以分为满秩（规范型对角线上没有$0$）和不满秩（规范型对角线上有$0$）两种.

\medskip

这样我们就对广义内积结构进行了完全的分类.
这个分类既可以从广义内积本身叙述，
也可以从度规矩阵（实对称矩阵）和二次型的角度来叙述.
当然最简单的是从规范型的角度来叙述（对角线上是否有$1,-1,0$）.
这些叙述方式的切换需要熟记于心.

最后做两点补充：

(1)~一个实对称矩阵的秩和正负惯性指数从合同标准型也能看出来.
而相似标准型（对角线为$n$个实特征值）就是一个合同标准型.
因此，
\textbf{实对称矩阵的非零特征值的数量就是它的秩，
正特征值的数量就是正惯性指数.}
从解题的角度，
这一结论沟通了合同与相似这两种不同的关系，
并且可以由此衍生出一些结论.
比如：一个实对称矩阵是正定的，
当且仅当它的所有特征值都是正数.
有一些问题在合同的框架下看起来没法处理，
联系到相似关系就能轻松处理了.
比如：为什么$A$正定可以推出$A^{-1}$正定？


(2)~实对称矩阵正定还有一个等价条件，
即\textbf{所有顺序主子式都大于$0$}.
这个定理在正定性相关的计算题中常用（见习题部分）.
关于它的证明，
后面也会结合其他相关习题一起讲解.


\subsection{*复内积空间（暂不全）}

实内积是以$\mathbb{R}^n$上的标准内积为模板的，
那么复内积应当是以$\mathbb{C}^n$上的标准内积为模板.
因此我们首先要讲清楚$\mathbb{C}^n$上的标准内积如何定义.
我们当然希望他和$\mathbb{R}^n$中的标准内积一样：
设$\bm{\alpha} = (a_1 , a_2 , \cdots , a_n)^T$，
$\bm{\beta} = (b_1 , b_2 , \cdots , b_n)^T$
则它们的内积（希望）定义成：

\begin{align*}
    (\bm{\alpha} , \bm{\beta}) = a_1 b_1 + a_2 b_2 + \cdots + a_n b_n
\end{align*}

但不妙的是这种定义将不会满足正定性.
比如$(i,i) \cdot (i,i) = -2 ,~(i,1) \cdot (i,1) = 0$.
这是不可接受的.
我们希望内积满足正定性，
是因为我们希望进一步地定义向量的长度为自身内积的平方根.
一个向量与自身的内积若不是非负实数，
在直观上是没意义的（果真如此吗？）.
为此我们对标准内积的定义做一些调整，
牺牲一些别的性质来保护正定性.
方法是在做内积时让某一个向量取复共轭，
使得$(\bm{\alpha} , \bm{\alpha})$中辐角相互抵消.
于是，
$\mathbb{C}^n$中的标准内积定义为：

\begin{align*}
    (\bm{\alpha} , \bm{\beta}) = a_1 \overline{b_1} + a_2 \overline{b_2 }+ \cdots + a_n \overline{b_n}
\end{align*}

这样定义的复标准内积就满足了正定性，
但是牺牲了对称性.
容易验证，
取而代之的新性质是\textbf{共轭对称性}：
$(\bm{\alpha} , \bm{\beta}) = \overline{(\bm{\beta},\bm{\alpha})}$.

以$\mathbb{C}^n$中的标准内积为模板，
可以叙述一般复内积的定义：

\medskip

\textit{设$V$是$\mathbb{C}$上的线性空间.
定义$V$上的一个二元函数$(\cdot,\cdot) : V \times V \to \mathbb{C}$，即：
对于任意$\bm{x},\bm{y} \in V$，$(\bm{x},\bm{y})$都对应一个复数.
如果对任意的$\bm{x},\bm{y} \in V$，$\alpha \in \mathbb{C}$，
满足以下$4$条规则，
那么称$V$是复内积空间（酉空间），
这个二元函数是$V$上的一个复内积：}

(1)~\textbf{正定性：}\textit{$(\bm{x},\bm{x} )\ge 0$，等号成立当且仅当$\bm{x} = \bm{0}$.}

(2)~\textbf{共轭对称性：}\textit{$(\bm{x},\bm{y}) = \overline{(\bm{y},\bm{x})}$.}

(3)~\textbf{线性（数乘）：}\textit{$(\alpha \bm{x}, \bm{y}) = \alpha (\bm{x},\bm{y})$.}

(4)~\textbf{线性（加法）：}\textit{$(\bm{x}+\bm{y},\bm{z}) = (\bm{x},\bm{z}) + (\bm{y},\bm{z})$.}\vspace{1em}


\medskip


\medskip

复内积的表示与实内积的思路是类似的，
但因为复内积满足的是共轭对称性，
表达式上有一些调整.
首先，
对内积进行展开得到的是：
\begin{align*}
    (\bm{\alpha} , \bm{\beta}) = 
    \left( \sum_{i=1}^{n}a_i \bm{e}_i , \sum_{j=1}^{n}b_j \bm{e}_j \right)
    = \sum_{i=1}^{n} \sum_{j=1}^{n} a_i \overline{b_j} g_{ij} 
\end{align*}

写成矩阵形式是：
\begin{align*}
    (\bm{\alpha} , \bm{\beta})
    = (a_1 , a_2 , \cdots ,a_n) 
    \begin{pmatrix}
        g_{11} & g_{12} & \cdots & g_{1n} \\
        g_{21} & g_{22} & \cdots & g_{2n} \\
        \vdots & \vdots  & \ddots & \vdots \\
        g_{n1} & g_{n2} & \cdots & g_{nn}
    \end{pmatrix}
    \begin{pmatrix}
        \overline{b_1} \\
        \overline{b_2} \\
        \vdots \\
        \overline{b_n}
    \end{pmatrix}
\end{align*}

这里$g_{ij}$给出一个复矩阵$H$.
$\bm{\alpha}, \bm{\beta}$在给定基下的内积表示为$(\bm{\alpha} , \bm{\beta}) = \bm{X}^T H \overline{\bm{Y}}$.
由于复内积有共轭对称性，
$H$与自身的共轭转置$\overline{H}^T$相等.
我们把复矩阵$H$的共轭转置记为$H^{\dagger}$.
则$H$满足$H=H^{\dagger}$.
满足此一性质的复矩阵称为Hermite矩阵.
可以把共轭转置看成转置在复矩阵中的推广，
把Hermite矩阵看成实对称矩阵的推广.

由内积的正定性可知，
$\bm{X}^T H \overline{\bm{X}} > 0$对于任意非零向量$\bm{X}$成立，
因此复度规矩阵$H$也是正定矩阵.
表达式$\bm{X}^T H \overline{\bm{X}}$称为Hermite型.

由此看出，
\textbf{酉空间中的复内积与正定Hermite矩阵之间存在一一对应.}

\clearpage

\subsection{习题}

\subsubsection{合同标准型和规范型}


本章的典型计算题是把一个实对称矩阵化为合同标准型/规范型.
注意看清题目究竟是问标准型还是规范型.

化合同标准型的方法有三种：

(1)~直接配方法.
这种方法一般是最快的，
很适合做计算题时使用.
缺点在于它纯粹是代数运算，
理论上的意义不大.

(2)~成对初等行列变换法.
合同变换式$C^T A C$，
$C$可以看成对$A$进行了一系列初等列变换，
而$C^T$则是进行了一系列\textbf{相对应的}初等行变换.
因此，
可以在$A$下放并排摆放一个单位矩阵，
然后与$A$一同做成对的初等行列变换，
使得$A$化为标准型.
此时，下方的单位矩阵就记录着所做的所有列变换，
因此它经历变换后得到的就是$C$.
这个方法速度也不慢，
做计算题时可用，
而且做一些证明题时需要用到这个理解.

(3)~相似对角化.
因为实对称矩阵总是可以正交相似对角化，
而正交相似变换就是一种合同变换，
所以相似标准型就是一种合同标准型.
直接用上一章的办法去计算正交相似对角化即可.
这个方法比较慢，
对于计算题来说是下策，
但是做一些证明题时需要用到这个理解.

\bigskip

\begin{example}
    思考：
    
    (1)~对于$n$阶实对称矩阵，

    有多少个合同等价类？
    
    (2)~对于复对称矩阵而言，“惯性指数”是什么？
\end{example}

\bigskip

\begin{example}
    将下述二次型化为标准型，
    并写出所做的非退化线性替换.

    (1)~$f(x_1, x_2, x_3) = x_1^2 + x_2^2 + x_3^2 + x_1 x_2 + x_1 x_3 + x_2 x_3$.
    
    (2)~$f(x_1, x_2, x_3) = x_1 x_2 + x_1 x_3 + x_2 x_3$.
\end{example}

\bigskip

\begin{example}
    设实二次型
    \[
    f(X)=f(x_1,x_2,x_3)
    =x_1^2+x_2^2+x_3^2+2ax_1x_2+2ax_1x_3+2ax_2x_3,
    \quad a\in\mathbb{R}\text{为参数}
    \]
    可用非退化线性替换 \(X=CY\) 化成二次型
    \[
    g(Y)=g(y_1,y_2,y_3)
    =y_1^2+y_2^2+4y_3^2+2y_1y_2
    \]
\end{example}


（1）求参数 \(a\) 的值.

（2）求把 \(f(X)\) 化成 \(g(Y)\) 的非退化线性替换中的可逆矩阵 \(C\).

\subsubsection{正定二次型}

判断一个实对称矩阵是否正定，
常用的办法是验证顺序主子式是否全部大于$0$.



\begin{example}
    判断$I_n + J_n$是否为正定矩阵
\end{example}

\begin{remark}
    可以去计算顺序主子式，
    但是更快的方法是去证$n$个特征值都为正.
\end{remark}

\bigskip


\begin{example}
    ~
    
    (1)~若二次型
    \[
    f(x_1,x_2,x_3)=2x_1^2+x_2^2+x_3^2+2x_1x_2+t\,x_2x_3
    \]
    是正定的，求 \(t\) 的取值范围.

    (2)~
    当 \(a_1,a_2,\cdots,a_n\) 满足什么条件时，实二次型
    \[
    f(x_1,x_2,\cdots,x_n)
    =(x_1+a_1x_2)^2+(x_2+a_2x_3)^2+\cdots+(x_{n-1}+a_{n-1}x_n)^2+(x_n+a_nx_1)^2
    \]
    是正定的？
\end{example}

\begin{remark}
    第(1)问比较常规，
    可用顺序主子式来计算.
    但是第(2)问要这样做就有点难了.
    请注意第(2)问已经是一个半正定的二次型.
\end{remark}

\bigskip


\begin{example}
    设$A$是$\mathbb{R}$上的$n$阶正定矩阵，
    任取列向量$\bm{\alpha} \in \mathbb{R}^n$，
    有：
    \[
    \begin{vmatrix}
        A & \bm{\alpha} \\
        \bm{\alpha}^T & 0
    \end{vmatrix} < 0
    \]
\end{example}

\begin{remark}
    通过成对的分块初等行列变换，
    将其合同变换为分块对角矩阵.
\end{remark}

\bigskip

\begin{example}
    设矩阵$A$：
    \[
    \begin{pmatrix}
        3 & -1 & 1 \\
        -1 & 3 & 1 \\
        1 & 1 & 3
    \end{pmatrix}
    \]

    (1)~证明$A$是正定矩阵.

    (2)~求可逆正定矩阵$C$，
    使得$C^2 = A$.

    (3)~思考：设$A$是任意实对称矩阵，
    存在正定实对称矩阵$C$使得$C^2 = A$是$A$正定的充要条件吗？
\end{example}

\begin{remark}
    对于第(2)问和第(3)问，考虑如果$A$是标准型，
    能否写出符合题意的$C$？
    现在$A$本身不是一个标准型，
    如何利用合同变换联系到标准型的分解？
\end{remark}

\bigskip



\begin{example}~

    (1)~证明：$n$ 级实对称矩阵 $A$ 正定的充要条件是，存在 $m\times n$ 列满秩实矩阵 $P$，使得
\[
A = P^{T}P
\]


    (2)~证明：$n$ 级实对称矩阵 $A$ 半正定的充要条件是，存在 $r\times n$ 行满秩实矩阵 $Q$，使得
\[
A = Q^{T}Q
\]
\end{example}

\begin{remark}
    考虑如果$A$是标准型，
    能否做出符合题意的分解？
\end{remark}


\subsubsection{数学归纳法}

在二次型问题中，
对矩阵阶数$n$做数学归纳法来证明一些结论.
关键的操作是：

(1)~把$n$阶矩阵$A$切出一行一列来，形成$2 \times 2$分块矩阵.

(2)~做合同变换（成对分块初等行列变换），使其化为分块对角矩阵.
在这个过程中，对于左上角的$(n-1) \times (n-1)$矩阵块，
可以利用归纳假设来处理.

\bigskip


\begin{example}
    证明：$n$阶实对称矩阵$A$正定的充要条件是$A$的所有顺序主子式都大于$0$.
\end{example}

（参考教材证明）

\bigskip

\begin{example}
    设 \(A\) 为实对称正定矩阵，试证明：必存在对角元都是正数的下三角矩阵 \(L\)，使得
    $A=LL^T$并且这个分解是唯一的.
\end{example}

\begin{proof}
    对矩阵$A$的阶数$n$做数学归纳法.

    (i)~当 $n=1$ 时，
    $A=(a)$，$a>0$，
    则 $(a)=(\sqrt{a})(\sqrt{a})$，结论显然成立.

    (ii)~
    当 $n>1$ 时，
    假设结论对任意 $n-1$ 阶实对称正定矩阵都成立.
    现有 $n$ 阶实对称正定矩阵 $A$.
    把 $A=(a_{ij})$ 分块成：
    \[
    A=
    \begin{pmatrix}
    A_{n-1} & \bm{\alpha} \\[2pt]
    \bm{\alpha}^T & a_{nn}
    \end{pmatrix},
    \]
    其中 $A_{n-1}$ 是$n-1$阶实对称矩阵.

    $A$ 是正定矩阵，
    则$A$的所有顺序主子式都为正，
    而$A_{n-1}$ 的顺序主子式就是 $A$ 的前$n-1$个顺序主子式，
    所以 $A_{n-1}$ 正定.

    对$A$做成对的分块初等行列变换
    （在下式中用初等分块矩阵乘法表示），
    化为分块对角矩阵：
    \[
    \begin{pmatrix}
    I_{n-1} & 0 \\[2pt]
    -\bm{\alpha}^T A_{n-1}^{-1} & 1
    \end{pmatrix}
    \begin{pmatrix}
    A_{n-1} & \bm{\alpha} \\[2pt]
    \bm{\alpha}^T & a_{nn}
    \end{pmatrix}
    \begin{pmatrix}
    I_{n-1} & -A_{n-1}^{-1}\bm{\alpha} \\[2pt]
    0 & 1
    \end{pmatrix}
    =
    \begin{pmatrix}
    A_{n-1} & 0 \\[2pt]
    0 & b
    \end{pmatrix}
    \]
    其中
    $b = a_{nn} - \bm{\alpha}^T A_{n-1}^{-1}\bm{\alpha}$.
    两边取行列式，得：
    \[
    |A| = |A_{n-1}|b
    \]

    因为$A$是正定矩阵，
    所以$|A|$和$|A_{n-1}|$都是正数，
    则$b$也是正数.
    将上式整理成对$A$的分解式：
    \[
    A=
    \begin{pmatrix}
    A_{n-1} & \alpha \\[2pt]
    \alpha^T & a_{nn}
    \end{pmatrix}
    =
    \begin{pmatrix}
    I_{n-1} & 0 \\[2pt]
    -\alpha^T A_{n-1}^{-1} & 1
    \end{pmatrix}^{-1}
    \begin{pmatrix}
    A_{n-1} & 0 \\[2pt]
    0 & b
    \end{pmatrix}
    \begin{pmatrix}
    I_{n-1} & -A_{n-1}^{-1}\alpha \\[2pt]
    0 & 1
    \end{pmatrix}^{-1}
    \]

    即：
    \[
    A = L_2
    \begin{pmatrix}
    A_{n-1} & 0 \\[2pt]
    0 & b
    \end{pmatrix}
    L_2^T,~~
    \text{其中下三角矩阵}
    L_2 = 
    \begin{pmatrix}
    I_{n-1} & 0 \\[2pt]
    -\alpha^T A_{n-1}^{-1} & 1
    \end{pmatrix}^{-1}
    \text{对角元为正数}
    \]
    

    由于 $b>0$，且 $A_{n-1}$ 正定，由归纳假设，
    存在对角元为正数的下三角矩阵$L_1$，
    使得$A_{n-1} = L_1 L_1^T$.
    另外$b = \sqrt{b}\cdot \sqrt{b}$.
    于是有：
    \[
    \begin{pmatrix}
    A_{n-1} & 0 \\[2pt]
    0 & b
    \end{pmatrix}
    =
    \begin{pmatrix}
    L_1 & 0 \\[2pt]
    0 & \sqrt{b}
    \end{pmatrix}
    \begin{pmatrix}
    L_1 & 0 \\[2pt]
    0 & \sqrt{b}
    \end{pmatrix}^T.
    \]

    因此，$A = L_2 L_3 L_3^T L_2^T$.
    记$L = L_2 L_3$，
    则$L$是对角元为正数的下三角矩阵，
    且$A = LL^T$.
    因此结论对$n$阶矩阵$A$也成立.

\end{proof}

\bigskip

\begin{example}
    已知$n$阶实对称矩阵$A$的各阶顺序主子式$\Delta_1 , \Delta_2 , \cdot , \Delta_n$
    都不为零.
    证明：
    二次型$f(\bm{X}) = \bm{X}^T A \bm{X}$可由非退化线性替换$\bm{X} = C \bm{Y}$化为下述标准型：
    \[
    f = \lambda_1 y_1^2 + \lambda_2 y_2^2 + \cdots + \lambda_n y_n^2
    \]

    其中$\lambda_k = \frac{\Delta_k}{\Delta_{k-1}}$，$\Delta_0 = 1$.
\end{example}





\subsubsection{*二次型的应用}

\begin{example}
    设二元实值函数$F(x,y)$有一个稳定点$\bm{\alpha} = (x_0 , y_0)$.
    $F(x,y)$在$(x_0 , y_0)$邻域内有三阶连续偏导数.
    如何利用$F(x,y)$在$(x_0 , y_0)$处的各个二阶偏导数判断该点是否为极值点？
    （回忆高等数学中的定理）
    相关结论能否推广至$n$元函数？
\end{example}

\bigskip

\begin{example}
    ~

    (1)~$\mathbb{R}^3$中的
    曲面$z = 5x^2 + 4xy + 2y^2$是什么类型的二次曲面？

    (2)~$\mathbb{R}^2$中的曲线$5x^2 + 4xy + 2y^2 - 24x - 12y + 18 = 0$是什么类型的二次曲线？
\end{example}